
The popularity paradox is real: models trained on high-popularity datasets exhibit generalization gaps 21.21 percentage points larger than those trained on low-popularity alternatives. Popular benchmarks attract 7.56 times more optimization research, consistent with a co-evolution hypothesis. However, texture bias is not the mechanism---modern architectures show improved shape-based generalization compared to legacy architectures.

The findings support complementing popular benchmarks with diverse alternatives. The benchmarks most trusted by the community may warrant the most scrutiny before deployment.

\subsubsection{Future Work}

\begin{itemize}
\item Test alternative artifact mechanisms (spurious correlations, frequency-domain biases)
\item Extend to multiple domains (NLP, tabular data)
\item Replicate at ImageNet scale
\item Conduct multi-run experiments for formal statistical testing
\end{itemize}

